\section{Introduction}
\label{sec:intro}

A multi-agent LLM system shares state. One agent writes a plan that another
reads; a tool registry is consulted at planning time and called at execution
time; a supervisor retracts a decision that a worker has already acted on. Each
agent's step has the shape of a transaction---read, compute, commit---with one
difference from every setting in which transactions have been studied: the
compute phase is a neural-inference call that takes seconds to minutes, during
which the read set is neither locked nor cheaply re-validated. Deployed
frameworks place this phase inside a pinned snapshot (LangGraph's superstep), a
re-rendered context (Letta) or a message queue (AutoGen), and none of them
states what guarantee the shared state carries across it.

\subsection{A failure, constructed and then observed}
\label{sec:intro-example}
Two agents share a trip record. A flight agent reads ``trip date $=$ 14 June,''
drafts a reservation over a model call, and books it. During that call a second
agent, acting for the user, moves the trip to 21 June. The flight agent commits
a 14 June booking beside a 21 June trip date. No component failed; the history
is even serializable, in the order in which the flight agent runs to completion
and the revision follows. What it is not is serializable \emph{in the order in
which its operations committed}: the revision committed first, and a flight
agent executing after it would have read 21 June. \Cref{fig:window} draws the
shape.

We ran this scenario on an unmodified LangGraph runtime under live models
(\cref{sec:live-a1}). With the second agent started inside the first agent's
generation window, the framework committed the stale booking in every session
in which the revision landed inside the window---$\LiveGptVanillaAOne\%$ and
$\LiveHaikuVanillaAOne\%$ of sessions on two API models, the remainder landing
after the commit---and read-set validation at commit removed every such firing.
The same data record a second fact the predicate does not see: when the
revision lands \emph{after} the commit, the booking contradicts the final trip
date just the same, and no isolation discipline can prevent it. We report both
numbers throughout, because the distance between them is the distance between a
consistency predicate and the harm it is meant to track. The rates themselves
are properties of the race we opened (a fixed $150$\,ms offset) and we attach
no prevalence reading to them.

\smallskip\noindent\textbf{What this paper is not claiming.} A census of
\CensusRepos{} public LangGraph repositories (\cref{sec:prevalence}) finds
the structural precondition of the failure above---a reader and a writer of
one last-value channel scheduled together---in none of \CensusSites{} graph
sites. We take that at face value: public LangGraph code is mostly sequential,
and four of five frameworks we probed close the window by construction. The
regime studied here is the one in which contention is designed in (parallel
fan-out over a shared store, several agents acting for one user, a registry
mutated while plans are in flight), and the contribution is not that the
anomalies are common: it is a catalog with mechanized definitions,
disciplines priced per unit of sound work, and an account of what each
leaves behind.

\begin{figure}[t]
\centering
\begin{tikzpicture}[x=0.62cm,y=0.52cm,font=\scriptsize,>=Latex]
  % time axis
  \draw[->] (0,0) -- (9.0,0) node[right]{$t$};
  \foreach \x/\l in {0.4/0,2.4/1,5.2/2,8.0/3} {\draw (\x,0.08)--(\x,-0.08) node[below]{\l};}
  % flight agent
  \node[anchor=east,align=right] at (-0.1,2.2) {flight\\agent};
  \draw[fill=black!8,draw=black!60] (0.4,1.8) rectangle (8.0,2.6);
  \node at (4.2,2.2) {generation phase (inference)};
  \draw[->,thick] (0.4,1.8) -- (0.4,0.3) node[midway,left]{read $c$};
  \node[anchor=south] at (0.4,2.6) {$\rt=0$};
  \draw[->,thick] (8.0,1.8) -- (8.0,0.3) node[midway,right]{commit $c'$};
  \node[anchor=south] at (8.0,2.6) {$\wt=3$};
  % date agent
  \node[anchor=east,align=right] at (-0.1,1.0) {date\\agent};
  \draw[fill=black!20,draw=black!60] (2.4,0.7) rectangle (5.2,1.3);
  \draw[->,thick] (5.2,0.7) -- (5.2,0.3) node[midway,right]{write $c$};
  \node[anchor=south] at (3.8,1.3) {$\wt=2$};
  % witness brace
  \draw[decorate,decoration={brace,amplitude=4pt,mirror}] (0.4,-0.75) -- (8.0,-0.75)
    node[midway,below=4pt,font=\scriptsize,align=center]{$\rt_i<\wt_j<\wt_i$,\ $\rv_i[c]\neq\wv_j[c]$:\\an \A1 witness};
\end{tikzpicture}
\caption{The read--generate--commit window. Logical time is the length of the
operation log. Operation $i$ (flight agent) reads $c$ at $\rt_i=0$; operation
$j$ (date agent) commits a different value to $c$ at $\wt_j=2$, inside $i$'s
window; $i$ commits at $\wt_i=3$ grounded in the superseded read. The
$L_1$ discipline re-validates the read set at $\wt_i$ and discards the
generation.}
\label{fig:window}
\end{figure}

\subsection{Why existing theory does not operationalize these failures}
\label{sec:why-classical-fails}
Three traditions bear on the problem and none applies as it stands. Hardware
models~\cite{Lamport1979SC,GharachorlooEtAl1990,PulteEtAl2018ARMv8,HerlihyWing1990}
presuppose bounded latency and static access sets; database isolation
theory~\cite{BerensonEtAl1995,Adya1999,AdyaLiskovOneil2000} supplies the
anomaly-by-anomaly method we follow but holds the read set under lock or
snapshot for the transaction's lifetime; replicated-storage
consistency~\cite{ViottiVukolic2016Consistency,LloydEtAl2011COPS,Burckhardt2014,TerryEtAl1994}
treats operations as discrete events with summarizable signatures. The agent
regime has unbounded generation latency, a tool registry that is first-class
mutable state, and irreversible external effects---so the cascading-abort
machinery of database recovery, which assumes effects can be undone, is
exactly what fails.

\subsection{Contributions}
\label{sec:contributions}
We claim no new concurrency theory. Each anomaly below is a classical one,
and we say which. The contribution is the operational port of four of them to
agent runtimes, its mechanization at stated assurance levels, and an
evaluation that measures the gap between each predicate and the harm it
tracks---including the part of the harm that is not a concurrency anomaly at
all.

\smallskip\noindent\textbf{(C1) A catalog with classical names, and the
residue it leaves} (\cref{sec:anomalies}). Four history predicates in
\TLA{}---stale generation (\A1), phantom tool (\A2), causal cascade (\A3),
tool-effect reordering (\A6)---each with a TLC witness and its classical
counterpart: a failed backward validation~\cite{KungRobinson1981OCC}, a
dangling capability reference, recoverability with externalization as the
commit point~\cite{BernsteinHadzilacosGoodman1987}, and order preservation
of external effects. \Cref{thm:corc} (TLAPS, \TlapsCommitOrder{}
obligations) places \A1-freedom among the classical classes exactly, as a
conditional guarantee whose hypotheses (per-agent non-overlap,
read-before-commit, faithful records) are stated, checked on every committed
record, and violated by one deployed construct we name. Beside the catalog
we define the \emph{residue predicate} \Rp{}: an irreversible effect is out
while ordinary later writes have revised the state it was derived from. \Rp{}
is not an anomaly---it fires in correct sequential executions---and we say so;
it is the harm no commit-time discipline reaches, decided on the records by
argument provenance, and on the live study \A1 and \Rp{} partition the
outcome column (\cref{sec:residue}).

\smallskip\noindent\textbf{(C2) A lattice, realized} (\cref{sec:lattice}). The
$2^4$ exclusion lattice is trivial; what is not is that every point is inhabited
by a checked runtime configuration (\LatticeRuns{} TLC runs, no mismatch,
bounded as stated), that one chain $L_0\subsetneq\dots\subsetneq L_4$ has a
verified realizing runtime at every level, and that the basis ledger that
answers \Rp{} composes with every point without disturbing its verdict
(\LatticeThirtyTwoRuns{} further runs).

\smallskip\noindent\textbf{(C3) Verified detectors and disciplines}
(\cref{sec:mechanization}). \VerusCurated{} Verus obligations, no
\texttt{assume} or \texttt{admit}, under two enumerated axioms and three
enumerated stubs: detectors sound and complete against their specifications,
two concurrent $L_1$ stores proved window-free under a verified reader--writer
lock, $L_2$ (output commit), $L_3$ (commit-order sequencing) and $L_4$
(validation at dispatch) runtimes whose $L_2$ and $L_4$ bytes are the ones
measured. \Cref{tab:assurance} labels every claim with its assurance
level---checked model, verified decision procedure, verified bytes, unverified
port, open---and names each correspondence between a \TLA{} definition, its
Verus restatement, the executed code and the scorer, with the hypothesis
under which it holds and the gate that checks the hypothesis on the data.
We account for the obligations by what they bear (supplement \S\,S):
\VerusLoadBearing{} underpin a deployed or measured artifact and
\VerusScaffold{} are scaffolding. The basis ledger that answers \Rp{} is
\LfiveVerusIntro{} in Verus: its invariant is preserved by every step and
excludes the residue witness; its executable form in the live study is
harness code, labeled as such.

\smallskip\noindent\textbf{(C4) An evaluation scored on outcomes, not flags}
(\cref{sec:eval}). For each discipline we report the anomaly it excludes, the
harm that remains, and the price: on the live LangGraph study, the $L_1$
residue is \Rp{} in $\LiveGptValidatedWrong\%$ and
$\LiveHaikuValidatedWrong\%$ of sessions at the $150$\,ms offset, decided on
the records; with the commit decision taken by the verified $L_1$ store
through a $100$-line bridge the residue is the same kind of number
(\LiveVerifiedOutcome), and the basis ledger with a compensator removes it
(\LfiveLiveOutcome) at one compensation generation per supersession; on
$1{,}000$ seeded cascades per
depth, output commit removes exposure exactly when the review precedes the
release and refuses the retraction otherwise; per sound committed operation,
snapshot isolation costs $+0.4$--$11.1\%$ and pessimistic locking
$+1.2$--$97\%$ in tokens, with a wall-clock measurement whose run-order
confound we quantify rather than assume away.

\smallskip\noindent\textbf{(C5) Where the preconditions occur} (\cref{sec:prevalence}).
A census of \CensusRepos{} public LangGraph repositories (\CensusSites{} graph
sites) finds no co-scheduled reader and writer of a last-value channel---a
negative result for \A1's structural precondition, which we report as one---while
a released tool executor dispatches a model-emitted, order-dependent pipeline
concurrently on $\ToolNodeGptBatched/\ToolNodeGptBatched$ turns of one backbone
and never on another, which is the case for \A6.

\smallskip\noindent\textbf{What this paper does not establish.} It does not
establish a production prevalence for any anomaly: every occurrence figure is
a race our harness opened or a rate on our own workloads, and the census
measures how public code is written, not how often it runs. It does not verify
the instrumentation that produces the traces, nor the Python harnesses; the
trust base is enumerated in \cref{sec:trust}. And it does not claim that a
verified arm's zero is evidence: a verified arm's zero is a theorem executing,
so the evaluation leads with costs and with the counts that read no runtime
flag.
